\begin{helper}
Let $\mathcal F$ be a $k$-uniform hypergraph, let $f$ be an edge, and suppose
every vertex of $f$ has degree $D$.  Suppose also that every edge other than
$f$ which meets $f$ meets it in exactly one vertex.  For $v\in f$ put
\[
C_v=\{g\in E(\mathcal F):g\ne f,\ v\in g\},
\]
and let $N$ be the set of edges other than $f$ which meet $f$.  Then
$|C_v|=D-1$ for every $v\in f$, the sets $C_u$ and $C_v$ are disjoint for
distinct $u,v\in f$, every edge in $N$ lies in one of the classes $C_v$, and
the number of two-element sets of edges which lie inside a single class
$C_v$ is at most
\[
k\binom{D-1}{2}.
\]
\end{helper}
\begin{proof}
Fix $v\in f$.  By the definition of degree in
\noderef{HypergraphDegree}, the hypothesis that the degree of $v$ is $D$
means exactly that $D$ edges of $\mathcal F$ contain $v$.  One of them is
$f$ itself.  Removing $f$ from this degree count leaves exactly $D-1$ edges
other than $f$ containing $v$, namely the members of $C_v$.  Hence
$|C_v|=D-1$.

Now let $u,v\in f$ be distinct.  If some edge $g$ belonged to both $C_u$ and
$C_v$, then $g\ne f$ and $g$ would contain both $u$ and $v$.  Thus
$g\cap f$ would contain the two distinct vertices $u$ and $v$, contradicting
the hypothesis that every edge other than $f$ which meets $f$ meets it in
exactly one vertex.  Therefore $C_u$ and $C_v$ are disjoint.

If $g\in N$, then $g\ne f$ and $g$ meets $f$, so choose
$v\in g\cap f$.  Then $v\in f$ and $g\in C_v$.  Conversely, if
$g\in C_v$ for some $v\in f$, then $g\ne f$ and $v\in g\cap f$, so
$g\in N$.  Thus the classes $C_v$ cover exactly the neighbor set $N$.

For a fixed $v\in f$, the two-element sets of edges contained in $C_v$ are
precisely the two-subsets of $C_v$, and there are
$\binom{|C_v|}{2}=\binom{D-1}{2}$ of them.  Since $\mathcal F$ is
$k$-uniform, \noderef{UniformHypergraph} gives $|f|=k$.  The collection of
all two-element sets which lie inside at least one class is contained in the
union, over the $k$ vertices $v\in f$, of these two-subset collections.  A
union has cardinality at most the sum of the cardinalities of its members, so
the total number is at most
\[
\sum_{v\in f}\binom{D-1}{2}=k\binom{D-1}{2},
\]
as claimed.
\end{proof}
